% ===================================================================
%  도표 제 13 호 — 호텔. 식당. 찻집.
%
%  Traduction, faite en 2026, en coreen d'aujourd'hui, sur l'ido de
%  texto/io. Regles de la colonne : voir 10-dopyo-01.tex. Recit a la
%  premiere personne, sans scenes ; les cles ne portent ni c1 ni c2.
%  Les deux notes d'auteur sont en gras-italique dans le fac-simile,
%  contrairement a celles des tableaux 1, 5 et 12 : la colonne suit.
%
%  TRENTE-TROIS PAIRES POUR QUATRE-VINGT-TREIZE RENVOIS — le taux le
%  plus eleve de la colonne coreenne apres le tableau 6. La colonne
%  telougoue avait mesure exactement la meme chose au meme tableau, et
%  la colonne marathe avant elle. Le 13 ATTACHE : le conducteur DE
%  l'omnibus, la malle DU voyageur, la cuisine AU sous-sol, l'enseigne
%  DE la halle, la caissiere QUI inscrivait la facture SUR le
%  registre. Trois langues de trois familles, trois fois le meme
%  tableau en tete, et la cause est dans le texte de 1926.
%
%  LE MENU RESTE FRANCAIS, ET C'EST LA REGLE. Les tripes a la mode de
%  Caen, le gigot de mouton, le Saint-Emilion sont les CHOSES qu'on
%  mange dans cette salle-la ; les remplacer par des plats coreens
%  aurait deplace le repas au lieu de traduire la langue. La colonne
%  nomme donc ce qu'elle voit — 캉풍 곱창 요리, 양 넓적다리 구이,
%  크림 시금치 — comme le hindi et le telougou avant elle.
%
%  L'HOTELLERIE EST VIEILLE EN COREE ET TOUT S'Y DIT EN COREEN :
%  문지기 le portier, 요리사 le cuisinier, 하녀 la servante, 짐꾼 le
%  porteur, 빨래하는 여자 la blanchisseuse, 우편함 la boite aux
%  lettres, 잉크 l'encre, 봉투 l'enveloppe, 채찍 le fouet, 마부석 le
%  siege du cocher. L'automobile est jeune, et tout y reste anglais :
%  운전사, 정비공, 자동차, 엘리베이터, 자전거 타는 사람, 당구. Meme
%  coupure a l'annee qu'au tableau 13 telougou, sur les memes objets.
%
%  UN JEU QUE JE NE SAIS PAS NOMMER, comme au telougou : les « dami »
%  (79), les dames. Le coreen a ses jeux de plateau — 바둑, 장기,
%  윷놀이 — et aucun n'est le jeu de dames. La colonne ecrit
%  체커, l'emprunt courant, plutot que de coller le nom d'un jeu
%  coreen sur un autre jeu. Meme raison qu'au jeu de barres du
%  tableau 2.
% ===================================================================

%%K t13-noto-1 noto
\VUnotes{143.2mm}{%
\textit{\VUgras{(*) 침실의 자세한 것은 도표 제 6 호를 보라.}}%
}

%%K t13-apar-1 apar
\VUsaut{3.0mm}
\VUtitre{15.0pt}{60}{도표 제 13 호}
\VUsaut{1.6mm}
\VUpk{10.2pt}{40}{호텔. --- 식당.}
\VUsaut{2.0mm}
\VUpk{10.2pt}{40}{찻집.}
\VUsaut{3.4mm}

%%K t13-01-1 p
1. --- 어제 저녁 루아양에 닿아 나는 「\textit{오세앙 호텔}」에 들었다. 한 벗이 나에게
일러 준 곳이다.

%%K t13-01-2 p
사람들이 나에게 고운 \VUgras{방} \textsuperscript{(2)}을 내주었다 — 이 \VUgras{층}
\textsuperscript{(1)}에 있는 방이다 —, 거기에서 나는 아주 잘 잤다 (*).

%%K t13-02-1 p
2. --- 오늘 아침 내 방에서 나왔다 — 그 방으로 \VUgras{하녀} \textsuperscript{(3)}가
아침상을 들여왔었다 —, 그리고 \VUgras{담배 피우는 방} \textsuperscript{(5)}으로 들어갔다.
그 방은 \VUgras{읽는 방} \textsuperscript{(4)}으로도 쓰인다. \VUgras{신문}
\textsuperscript{(6)}을 읽으려는 것인데, \VUgras{담배} \textsuperscript{(8)}를 피우면서
읽었다. 읽고 나서 \VUgras{엘리베이터} \textsuperscript{(9)}로 내려가 시내를 거닐러 갔다.

%%K t13-03-1 p
3. --- 호텔 \VUgras{사무원} \textsuperscript{(12)}에게 묻는 것을 잊었다 —
\VUgras{공동 식탁} \textsuperscript{(11)}에서 몇 시에 밥을 내는지. 그래서 일찍 돌아와, 그
틈에 호텔 \VUgras{목욕실} \textsuperscript{(10)}에서 목욕을 했다. 그 다음에 벗 하나에게
전화하려고 다시 \VUgras{계산대} \textsuperscript{(15)}로 갔다. 그런데 \VUgras{전화기}
\textsuperscript{(13)}가 쓰이는 중이었다. 내 난처함을 보고 \VUgras{여자 회계}
\textsuperscript{(14)}가 — 그 여자는 그때 \VUgras{계산서} \textsuperscript{(17)}를
\VUgras{손님 명부} \textsuperscript{(16)}에 적고 있었다 — \VUgras{문지기} \textsuperscript{(18)}
를 불러 \VUgras{사환} \textsuperscript{(19)}을 보내 달라고, \VUgras{자전거로}
\textsuperscript{(20)} 내 심부름을 해 달라고 부탁해 주었다.

%%K t13-04-1 p
4. --- 나는 그 여자에게 고맙다고 하고, \VUgras{짐꾼} \textsuperscript{(24)}(품삯일을 하는
사람)에게 얼마쯤 주려고 나갔다. 그는 정거장에서 제 \VUgras{손수레} \textsuperscript{(25)}로
내 \VUgras{모자 상자} \textsuperscript{(26)}, \VUgras{내 트렁크} \textsuperscript{(27)},
\VUgras{내 여행 가방} \textsuperscript{(28)}을 날라다 주었다. 막 도착한 \VUgras{우편 배달부}
\textsuperscript{(21)}가 나에게 \VUgras{편지} \textsuperscript{(22)} 한 통을 주었다.

%%K t13-05-1 p
5. --- \VUgras{우편함} \textsuperscript{(23)} 곁, 문지방 위에서 나는 얼마 동안 더 서서
길의 오감을 바라보았다. \VUgras{마부} \textsuperscript{(30)}가 — \VUgras{합승 마차}
\textsuperscript{(29)}의 마부다 — 제 수레 위에 큰 \VUgras{궤} \textsuperscript{(31)}를
싣고 있었다. 그것은 어느 \VUgras{나그네} \textsuperscript{(32)}의 것이고, 그 나그네를
\VUgras{호텔 주인} \textsuperscript{(33)}이 배웅하고 있었다. 젊은 \VUgras{전보 배달부}
\textsuperscript{(35)}가 느긋하게 \VUgras{전보} \textsuperscript{(34)}를 호텔 손님 하나에게
가져왔다. \VUgras{나그네} \textsuperscript{(36)} 하나는 — 어깨에 \VUgras{사진기}
\textsuperscript{(38)}를 메고 — 제 \VUgras{안내 책} \textsuperscript{(37)}을 들여다보고
있었다.

%%K t13-tit-1 sub
\VUsaut{4.0mm}
\VUpk{10.2pt}{40}{점심때.}
\VUsaut{3.0mm}

%%K t13-06-1 p
6. --- 쇠울 앞에서 걱정 없는 \VUgras{심부름꾼} \textsuperscript{(39)} 하나가 제
\VUgras{짐 갈고리} \textsuperscript{(40)}와 제 \VUgras{구두약 상자} \textsuperscript{(41)}
사이에서 졸고 있었다. 그동안 젊은 \VUgras{빨래하는 여자} \textsuperscript{(42)}가 — 그
여자는 빨랫감 \VUgras{바구니} \textsuperscript{(43)}를 이고 갔다 — 문 곁에 서 있는
\VUgras{식당 우두머리} \textsuperscript{(44)}의 엄숙한 얼굴을 보고 웃었다.

%%K t13-07-1 p
7. --- \VUgras{요리사} \textsuperscript{(45)}가 보였다. 그는 제 \VUgras{부엌}
\textsuperscript{(47)}에서 나왔는데, 그 부엌은 \VUgras{지하층} \textsuperscript{(48)}에
있다. \VUgras{부엌 사환} \textsuperscript{(46)}을 심부름 보내려는 것이었다. 그때 점심때가
가까워진 것이 생각나 나는 식당으로 들어갔다. 가는 길에 안마당을 지났다. 거기에서
\VUgras{운전사} \textsuperscript{(50)}가 제 \VUgras{자동차} \textsuperscript{(49)}를
세워 두었고, \VUgras{정비공} \textsuperscript{(52)} 하나가 \VUgras{자전거 타는 사람}
\textsuperscript{(53)}이 일러 주는 대로, 방금 사고가 난 \VUgras{석유 세발차}
\textsuperscript{(51)}를 고치고 있었다.

%%K t13-08-1 p
8. --- 나는 젊은 \VUgras{사환} \textsuperscript{(54)} 하나에게 물었다 — 그는
\VUgras{맥주잔} \textsuperscript{(55)}을 나르고 있었다 — 공동 식탁이 베란다 아래에 있느냐고.

%%K t13-08-1 p suite
그가 그렇다고 하기에 나는 식탁 하나에 자리를 잡고 아주 훌륭한 밥상을 받았다.

%%K t13-noto-2 noto
\VUnotes{143.2mm}{%
\textit{\VUgras{(*) 낱말집에 실린 요리 목록의 도움으로, 선생님은 아래의 차림표를
쉽게 바꿀 수 있을 것이다.}}%
}

%%K t13-tit-2 sub
\VUsaut{4.0mm}
\VUpk{10.2pt}{40}{훌륭한 점심.}
\VUsaut{3.0mm}

%%K t13-09-1 p
9. --- 종업원이 나에게 점심 \VUgras{차림표} (*)와 포도주 목록을 가져왔다. 나는 골라서
\VUgras{곁들이}로 \VUgras{버터}와 \VUgras{새우}를, 첫 접시로 \VUgras{캉풍 곱창 요리}를
시켰다. \VUgras{생선}은 먹지 않고 양 \VUgras{넓적다리} 한 몫을 청했으며,
\VUgras{남새 요리}로 크림을 곁들인 \VUgras{시금치}를 시켰다. 그 다음
\VUgras{중간 요리} 가운데 마음에 드는 것이 없어서 여러 가지 후식을 가져오게 했다.
\VUgras{치즈}, \VUgras{과일}, \VUgras{비스킷}이다. 거기에 훌륭한 생테밀리옹
\VUgras{포도주}를 곁들였다. 밥상에 아주 흡족해서 나는 \VUgras{종업원} \textsuperscript{(57)}
에게 넉넉한 \VUgras{행하} \textsuperscript{(59)}를 주고 자리를 떴다.

%%K t13-10-1 p
10. --- 내가 떠날 채비를 하는 것을 보고 \VUgras{지배인} \textsuperscript{(60)}이 발코니로
가 보라고 권했다 — \VUgras{큰 바다} \textsuperscript{(62)}의 훌륭한 경치를 보라는 것이다.
멀리 \VUgras{고깃배} \textsuperscript{(63)}와 \VUgras{돛단배} \textsuperscript{(64)}와 거대한
\VUgras{여객선} \textsuperscript{(65)}이 보이기 시작했다. 그 배의 검은 그림자가 흰
\VUgras{구름} \textsuperscript{(67)} 위로 도드라진다 — \VUgras{수평선} \textsuperscript{(66)}
의 구름이다. 가벼운 바람이 \VUgras{깃발} \textsuperscript{(70)}을 흔들었다. 그것은
\VUgras{간판} \textsuperscript{(69)} 위에 꽂혀 있다 — \VUgras{장터 건물} \textsuperscript{(68)}
의 간판이다.

%%K t13-tit-3 sub
\VUsaut{3.0mm}
\VUpk{10.2pt}{40}{놀이.}
\VUsaut{3.0mm}

%%K t13-11-1 p
11. --- 그 경치가 어찌나 볼 만한지, 나는 내일 \VUgras{쌍안경} \textsuperscript{(71)}이나
\VUgras{망원경} \textsuperscript{(72)}을 들고 찻집 옥상에 올라가기로 마음먹었다.

%%K t13-12-1 p
12. --- 발코니를 떠나 나는 호텔 찻집으로 갔다. \VUgras{마실 것} \textsuperscript{(73)}을
한잔 하고 \VUgras{당구 한 판} \textsuperscript{(75)}을 치려는 것이다. 나는 멈추지 않고
찻방을 지났다. 거기에서는 많은 손님이 \VUgras{장기} \textsuperscript{(78)},
\VUgras{체커} \textsuperscript{(79)}, \VUgras{도미노 패} \textsuperscript{(80)},
\VUgras{트릭트락} \textsuperscript{(81)}, \VUgras{화투} \textsuperscript{(82)}를 하고
있었다. 나는 곧바로 \VUgras{당구방} \textsuperscript{(74)}으로 올라갔다.

%%K t13-13-1 p
13. --- 판이 끝나자 나는 아래층으로 내려가, 편지를 쓰려고 종업원에게 \VUgras{봉투}
\textsuperscript{(84)}, \VUgras{종이} \textsuperscript{(83)}, \VUgras{우표} \textsuperscript{(85)},
\VUgras{펜} \textsuperscript{(87)}, \VUgras{잉크} \textsuperscript{(86)}를 청했다.

%%K t13-13-2 p
그 다음에 \VUgras{마부} \textsuperscript{(92)} 하나를 불렀다. 그의 \VUgras{마차}
\textsuperscript{(88)}가 찻집 앞에 서 있었다. 시간으로 셈하는지 한 번 타는 것으로 셈하는지
값을 물었고, 값이 맞자 그는 나에게 \VUgras{마차 문} \textsuperscript{(89)}을 열어 주고
앉게 했다. 그는 다시 \VUgras{마부석} \textsuperscript{(91)}에 올라 \VUgras{채찍}
\textsuperscript{(93)}으로 말을 빠르게 몰았고, 우리는 즐거운 나들이를 떠났다.

\VUsaut{6.0mm}
\VUfilet{22mm}
